\section{Tournament Details}
\label{app:tournament_details}

This appendix provides the full tournament methodology and results for the
head-to-head evaluation referenced in Section~\ref{sec:tournament}: the
topic-selection procedure, per-topic outcomes, score distributions, full
statistical test battery, side-asymmetry and judge-bias visualizations,
and a comparison to an earlier tournament version.

\subsection{Topic Selection}
\label{app:topic_selection}

The 17 tournament resolutions were selected from a corpus of 54 candidate
topics through a judge-bias probe. For each candidate, we queried three
judge models (Claude Opus~4.5, Gemini~2.5~Pro, and GPT-5.2) with the
resolution and asked them to rate it on a 0--10 scale, where 0 indicates
complete support for the negative position, 5 indicates a neutral topic
with no directional preference, and 10 indicates complete support for the
affirmative. For each topic we computed the mean rating across judges and
the distance from 5.0; we selected the 17 topics with the smallest distance
as the tournament set.

A sample of selected topics with their bias distances:
\begin{itemize}
    \item ``The ends can justify the means'' (dist: 0.17)
    \item ``Mathematical objects exist independently of human minds'' (dist: 0.17)
    \item ``Deep learning will lead to artificial general intelligence'' (dist: 0.17)
    \item ``Moral truths are objective facts independent of human beliefs'' (dist: 0.30)
    \item ``Compulsory voting strengthens democracy'' (dist: 0.80)
\end{itemize}

This approach identifies topics where the chosen judge models express no
strong directional preference; it does not eliminate all sources of bias.
In particular, topics where judges are systematically wrong in the same
direction would appear unbiased by this metric but would still be biased
in absolute terms. The probing script is available at
\texttt{scripts/ipda\_experiment/iter3/tournament/probe\_topic\_bias.py}.

\subsection{Per-Topic Results}

\begin{table}[h]
\centering
\small
\begin{tabular}{p{6cm}ccr}
\toprule
\textbf{Topic} & \textbf{Wins} & \textbf{Win\%} & \textbf{Mean $\Delta$} \\
\midrule
\multicolumn{4}{l}{\textit{Strong topics ($\geq$75\% win rate)}} \\
Moral truths are objective facts & 3/4 & 75\% & $+10.0$ \\
Economic development causes democratization & 3/4 & 75\% & $+9.6$ \\
Remote work reduces productivity & 3/4 & 75\% & $+8.4$ \\
Direct instruction $>$ inquiry-based learning & 3/4 & 75\% & $+7.8$ \\
Compulsory voting strengthens democracy & 3/4 & 75\% & $+4.2$ \\
Economic sanctions are immoral & 3/4 & 75\% & $+3.0$ \\
Many-Worlds $>$ Copenhagen interpretation & 3/4 & 75\% & $+2.2$ \\
\midrule
\multicolumn{4}{l}{\textit{Split topics (50\% win rate)}} \\
The ends can justify the means & 2/4 & 50\% & $+11.5$ \\
Mathematical objects exist independently & 2/4 & 50\% & $+12.0$ \\
Voters capable of democratic accountability & 2/4 & 50\% & $+8.0$ \\
Language acquisition is primarily innate & 2/4 & 50\% & $+2.4$ \\
Eliminate teacher tenure & 2/4 & 50\% & $-1.6$ \\
UBI should replace means-tested welfare & 2/4 & 50\% & $-4.6$ \\
\midrule
\multicolumn{4}{l}{\textit{Weak topics ($\leq$25\% win rate)}} \\
Deep learning will lead to AGI & 1/4 & 25\% & $-4.1$ \\
Federal $>$ unitary government & 1/4 & 25\% & $-1.7$ \\
Saturated fat increases CVD risk & 1/4 & 25\% & $-13.5$ \\
International institutions constrain states & 1/4 & 25\% & $-7.4$ \\
\bottomrule
\end{tabular}
\caption{Per-topic tournament results (pipeline judge). The trained model excels on
abstract/philosophical topics but struggles on empirical/institutional topics.}
\end{table}

\subsection{Score Distribution}

\begin{figure}[h]
\centering
\includegraphics[width=\columnwidth]{fig_score_distribution.pdf}
\caption{Distribution of score differentials (trained $-$ base) across
68 debates. The distribution is right-skewed with mean $+2.7$, indicating
the trained model tends to win by larger margins than it loses.}
\label{fig:score_distribution}
\end{figure}

\begin{figure*}[h]
\centering
\includegraphics[width=\textwidth]{fig_topic_winrate.pdf}
\caption{Per-topic trained model win rate (pipeline judge, 4 debates per
topic). Green indicates $\geq$75\% win rate, yellow 50\%, red $\leq$25\%.
Annotations show mean score differential. The trained model dominates
abstract topics but struggles on evidence-dense empirical claims.}
\label{fig:topic_winrate}
\end{figure*}

Pipeline judge mean score difference: $+2.72$ (SD $= 16.38$).
The distribution is bimodal: the trained model wins decisively when it wins
(mean winning margin $+14.2$) but also loses decisively on certain
topic-side combinations (mean losing margin $-11.8$).

\begin{table}[h]
\centering
\begin{tabular}{lcc}
\toprule
\textbf{Score Diff Range} & \textbf{Count} & \textbf{\%} \\
\midrule
Trained $+10$+ & 22 & 32.4\% \\
Trained $+5$ to $+10$ & 8 & 11.8\% \\
Trained $+0$ to $+5$ & 7 & 10.3\% \\
Base $+0$ to $+5$ & 11 & 16.2\% \\
Base $+5$ to $+10$ & 10 & 14.7\% \\
Base $+10$+ & 10 & 14.7\% \\
\bottomrule
\end{tabular}
\caption{Score difference distribution (pipeline judge).}
\end{table}

\subsection{Statistical Tests}

\begin{table}[h]
\centering
\begin{tabular}{lcc}
\toprule
\textbf{Test} & \textbf{Statistic} & \textbf{$p$-value} \\
\midrule
Pipeline binomial (37/68) & -- & 0.545 \\
Panel binomial (38/68) & -- & 0.396 \\
Pooled binomial (198/340) & -- & 0.003 \\
Paired $t$-test (score diff) & $t(67) = 1.37$ & 0.175 \\
Wilcoxon signed-rank & $W = 909.5$ & 0.107 \\
Cohen's $d$ (pipeline) & 0.166 & -- \\
Side difference (Welch's $t$) & $t = 3.32$ & 0.002 \\
NEG wins binomial (23/34) & -- & 0.058 \\
AFF wins binomial (14/34) & -- & 0.392 \\
\bottomrule
\end{tabular}
\caption{Full statistical test battery for tournament results.}
\end{table}

\subsection{Side Asymmetry Visualization}

Figure~\ref{fig:side_asymmetry} visualizes the side-specific performance
reported in Table~\ref{tab:side_asymmetry}. The asymmetry holds whether
measured by win rate (41.2\% AFF vs.\ 67.6\% NEG), mean score differential
($-3.4$ AFF vs.\ $+8.9$ NEG), or Cohen's $d$ by side (0.09 AFF vs.\ 0.48 NEG).
Crucially, the base model also exhibits the same directional pattern (32.4\%
AFF vs.\ 58.8\% NEG from its own side perspective): the gap is structural
to IPDA, and training amplified rather than created it. We interpret this
as evidence that reactive argumentation (NEG, which speaks second and can
respond to concrete AFF claims) is more amenable to preference optimization
than generative argumentation (AFF, which must build a case from scratch).

\begin{figure}[h]
\centering
\includegraphics[width=\columnwidth]{fig_side_asymmetry.pdf}
\caption{Trained model win rate by debate side (pipeline judge, 68 debates).
The 26.4 percentage-point gap between affirmative (41.2\%) and negative
(67.6\%) win rates reflects a structural IPDA side advantage amplified
by training.}
\label{fig:side_asymmetry}
\end{figure}

\subsection{Judge Bias Analysis}

\begin{table}[h]
\centering
\small
\begin{tabular}{lcccc}
\toprule
\textbf{Judge} & \textbf{Trained\%} & \textbf{As AFF} & \textbf{As NEG} & \textbf{Retries} \\
\midrule
Gemini 2.5 & 54.4\% & 64.7\% & 44.1\% & 19\% \\
GPT-5.2 & 61.8\% & 29.4\% & 94.1\% & 22\% \\
Opus 4.5 & 60.3\% & 50.0\% & 70.6\% & 29\% \\
Sonnet 4 & 60.3\% & 32.4\% & 88.2\% & 43\% \\
\bottomrule
\end{tabular}
\caption{Per-judge voting patterns. GPT and Sonnet show extreme NEG-side
bias (94\%, 88\%); Gemini exhibits the opposite (65\% AFF). Opposing
biases partially cancel in panel consensus (55.9\%).}
\label{tab:judge_bias}
\end{table}

\begin{figure*}[h]
\centering
\includegraphics[width=\textwidth]{fig_judge_bias.pdf}
\caption{Per-judge voting patterns by trained model's debate side. GPT and
Sonnet show extreme NEG-side bias (94\%, 88\%), while Gemini exhibits the
opposite pattern (65\% AFF). Opus is the most balanced judge.}
\label{fig:judge_bias}
\end{figure*}

\subsection{Tournament V3 Comparison}

An earlier tournament (V3, 20 debates) showed no significant difference
(trained wins 9/20 = 45\%, $p = 0.82$). V3 had two infrastructure issues
that corrupted the signal: (i) panel judge API failures---Gemini and GPT
API keys were expired during parts of the run, leaving some debates scored
only by Claude judges, and (ii) 7 of the 20 debates had at least one
\texttt{null} speech due to generation timeouts, creating judge decisions
based on incomplete transcripts. V4 resolved both issues: all 68 debates
were generated to completion with retries on timeout, and all four panel
judges were healthy throughout the run. The 34-fold difference in
statistical power between V3 and V4 (20 vs.\ 340 judge-debate votes) is
the main reason V4 detects the effect that V3 missed; the underlying
trained-model advantage was present in both runs but only became
statistically distinguishable at V4's sample size.
